\documentclass[12pt]{article}
\usepackage[utf8]{inputenc}
\usepackage[T1]{fontenc}
\usepackage{hyperref}
\title{Iranian Sign Language Language Resources for K-12 Curriculum Development}
\author{Compiled by Claude Sonnet 5.5 (claude-sonnet-5-5) for Andr\'as Kornai}
\date{}
\begin{document}
\maketitle

\section{Introduction}
Iranian Sign Language (ISO 639-3 \texttt{psc}), also called Persian or Farsi Sign Language, is the deaf community sign language of Iran. It is not standardized nationwide. Regional variation exists, for example the Kermanshah variety and its fingerspelling differences from Tehran. No reliable estimate of the number of users was found. A 2021 proposal suggests the name ``Esharani''.

\section{Primary Educational Resources}
A Tehran university group published the \emph{Dictionary of Farsi Sign Language: Standardized Signs} from 1984 (four editions), and an earlier 1980 volume came from the national welfare organization for the deaf. Deaf communities have criticized the dictionary for borrowed ASL signs. Current availability was not established. Educational policy has reportedly phased out sign language teaching in schools for the deaf.
\begin{itemize}
\item Learning of foreign languages by deaf and hard-of-hearing learners in Iran (Persian-language article): \url{https://www.ensani.ir/fa/article/download/626261}
\end{itemize}

\section{Linguistic and Reference Materials}
\begin{itemize}
\item Iranian Sign Language (overview): \url{https://en.wikipedia.org/wiki/Iranian_Sign_Language}
\item Reflections on the names of the sign language of Iran, and a proposal: \url{https://acadeafic.org/2021/02/11/sign-language-of-iran/}
\item Kermanshah variety of Iranian Sign Language and its differences from the Tehran variety (Persian-language article): \url{https://www.ensani.ir/fa/article/626292/%DA%AF%D9%88%D9%86%D9%87-%DA%A9%D8%B1%D9%85%D8%A7%D9%86%D8%B4%D8%A7%D9%87%DB%8C-%D8%B2%D8%A8%D8%A7%D9%86-%D8%A7%D8%B4%D8%A7%D8%B1%D9%87-%D8%A7%DB%8C%D8%B1%D8%A7%D9%86%DB%8C-%D9%88-%D8%AA%D9%81%D8%A7%D9%88%D8%AA-%D8%A2%D9%86-%D8%A8%D8%A7-%DA%AF%D9%88%D9%86%D9%87-%D8%AA%D9%87%D8%B1%D8%A7%D9%86%DB%8C}
\item Variations in the Baghcheban Manual Alphabet in Iranian Sign Language: \url{https://independent.academia.edu/ZGhari}
\item Research article on deaf learners and English teaching (Journal of Language Teaching and Learning): \url{https://jlts.um.ac.ir/article_46142.html?lang=en}
\end{itemize}

\section{Religious and Community Resources}
None found.

\section{Digital and Audio Resources}
No online video dictionary or corpus was found in the search results.

\section{Curriculum Development Context}
The main gaps are a community-approved, non-ASL-influenced lexicon, sign-language instruction in schools for the deaf, and teachers proficient in sign. A curriculum would need deaf-led documentation of regional varieties and a decision on standardization.
\end{document}